\chapter{Numerical underflow}\label{chapter:numerical_underflow}

This appendix chapter discusses the numerical underflow encountered in our implementation of the efficient ML decoder \cite{efficient_decoding_algorithm}.

A collection of logical error rate curves illustrating the problem is shown in figure \ref{fig:exmaple_observed_numerical_underflow}.
The characteristic and possible causes of this numerical underflow mechanism are examined in section \ref{sec:num_underflow_the_problem}. 

\begin{figure}[ht]
    \centering
    \includegraphics[width=0.5\linewidth]{numerical_underflow/ml_test_example.pdf}
    \caption{
            Effect of numerical underflow on the resulting logical error rate curves. 
            The data was obtained from simulation of a code-capacity noise model, decoded with the efficient ML decoder.
        }
    \label{fig:exmaple_observed_numerical_underflow}
\end{figure}

But before going into detail about the problem, we give a rough outline on the principles behind the efficient ML decoder in section \ref{sec:num_underflow_decoder}.
This will give an intuition where to look for the possible numerical underflow, as well as motivate the possible mitigation strategies.

In the last section of this chapter, section \ref{sec:num_underflow_solutions}, we introduce the error mitigation process used in this thesis, 
as well as further approaches to avoid numerical underflow in an implementation of this algorithm.

At this point, I would like to once again express my gratitude to Aron Marton for his continuous support throughout this project.
He provided valuable insights and shared his expertise during every stage of the work on the decoder. 
From developing an understanding of the fundamental concept of the decoder, to providing help and tips to in the search for possible causes of the observed logical error rate behavior,
and finally to exploring potential solutions.

Moreover, he made his own implementation of the decoder available to me. 
This not only greatly helped in the development of my own implementation 
but also enabled the localization of the numerical underflow, as I was able to reproduce the same error in his implementation. 

\section{The Decoder}\label{sec:num_underflow_decoder}

In this section, we briefly review the decoder, focusing only on the relevant aspects for understanding how numerical underflow arises.
We restrict our discussion to the algorithm for decoding independent bit- or phase-flip noise presented in \cite{efficient_decoding_algorithm}\footnote{
Note that the decoder corrects either bit-flip or phase-flip errors in a single decoding step. 
Consequently, decoding under independent bit- and phase-flip noise is achieved by applying the decoder twice.

Furthermore, there appears to be an error, or at least an inconsistency in the notation, 
in the description of the independent noise model in \cite{efficient_decoding_algorithm}. 
The author assumes independent bit- and phase-flip noise with probabilities $p_X=p_Z$ and subsequently states that $p_Y=p_X^2$. 
However, combining independent bit- and phase-flip channels with probabilities $p_X$ and $p_Z$ results in the effective error probabilities
\[
    \tilde{p}_X = p_X(1-p_Z), \qquad
    \tilde{p}_Z = p_Z(1-p_X), \qquad
    \tilde{p}_Y = p_X p_Z,
\]
where $\tilde{p}$ denotes the probabilities in the combined channel, while $p$ denotes the probabilities of the independent channels.


}, which is the decoder configuration used throughout this thesis 

The algorithm solves the ML decoding problem, see section \ref{ssec:general_ML_decoder}, by explicitly computing the probabilities of the cosets $\text{Pr}(E_\sigma \bar{P}\mathcal{S})$ corresponding to the measured syndrome $\sigma$.
Since  $X$- and $Z$-errors are decoded independently in this algorithm\footnote{
    Note once again that this decoder is only designed to be able to decode either $X$- or $Z$-errors.
} and we only consider a single logical encoded qubit, 
the logical operators defining the two cosets are either $\bar{P} \in \{\bar{I},\bar{X}\}$ or $\bar{P} \in \{\bar{I},\bar{Z}\}$.

In general, the probability of an individual cosets decreases as the code distance increases, since the number of possible syndrome growth exponentially.
For the surface code, the average cosets probability can be expressed as 
\begin{equation}
    \langle \text{Pr}(E_\sigma \bar{P} \mathcal{S}) \rangle = 2^{-m/2}= 2^{-d\cdot(d-1)},
    \label{eq:mean_prop_coset_in_distance}
\end{equation} 
where we consider either $X$- or $Z$-errors, consistent with the assumption of the decoder.
Consequently, we expect the probabilities of individual cosets to become quite small.
Note that we are using a average here an not a median, the median is highly dependent on the physical error rate, as will be discussed later. 

The efficient decoding algorithm consists of two stages. 
First, representative error $E_\sigma$ corresponding to the measured syndrome is generated.
This representative error is then passed to the main part of the decoder, 
which computed the probabilities of the associated cosets.

Broadly speaking, the decoder utilizes a mapping from the surface code to fermionic Gaussian states, 
which are subsequently mapped to matchgate quantum circuits, thereby enabling an efficient computation of the cosets probabilities on classical computers.
The complete derivation and algorithm of this mapping are presented in \cite{efficient_decoding_algorithm}.
For the discussion in this appendix, we point out, that this process included the computation of norm of Gaussian states $\Gamma \geq 0$ and its renormalization.

The algorithm used to generate the conventional representation error $E_\sigma$ is sufficiently short to be quoted in full. 
For the case of $X$-errors \cite{efficient_decoding_algorithm} states:
\begin{quotation}
``Let [$\sigma_u$] be the syndrome[bit] of a site stabilizer $A_u$. 
We choose the desired error [$E_\sigma$] by connecting each site $u$ with a non-zero syndrome [$\sigma_u$] to the left boundary by a horizontal string of $X$-errors 
and adding all such strings modulo two. 
Note that [$E_\sigma$] can be constructed in time $\mathcal{O}(n)$.'' 
\end{quotation}
Here, we adapted the notation to match the notation used throughout this thesis. 
We refer to this method to construct the representative error as the \emph{Bravyi} approach.


\section{The Problem}\label{sec:num_underflow_the_problem}

Before the rejection rate was implemented as a mitigation strategy, the logical error rate exhibited the form shown in \ref{fig:exmaple_observed_numerical_underflow}.
The simulation was performed using code-capacity noise model with only bit-flip errors, measuring the corresponding logical $Z$ observable.
The logical error rate $p_{log}$ generally follows the expected behavior for high physical error rates $p_{phy}$.
In particular the logical error rate curves for the code distances $d\in \{3,5,7\}$, show no obvious anomaly and they exhibit the asymptotic scaling predicted by equation \ref{eq:p_log_asymptotic_behavior}, as discussed in section \ref{ssec:threshold_theorem}.

Starting with the logical error rate curve of a distance $d=9$ surface code, an anomalous upward trend becomes apparent.
We can observe that the logical error rate exhibits an upward trend, 
breaking from the expected asymptotic scaling, at a physical error rates of $p_{phy} \approx 5\%$.
We will refer to this point as the \emph{critical point}.
Below the critical point, the logical error rate exhibits large fluctuations between different physical error rates, together with an overall slight upward trend as the physical noise rate decreases.
Considering the higher distances, we observe that the critical point is shifting to higher noise rates as the code distance is increased. 

Interestingly, the irregularity of the logical error rate curves below the critical point appears to not be random. 
Repeating the simulation yield the same curve, up to the expected statistical uncertainty. 
We therefore conclude that the observed anomaly is deterministic, rather than a statistical artifact.

As shown in equation \ref{eq:mean_prop_coset_in_distance}, the average probability of an individual coset decreases with increasing code distance.
Likewise, reducing the physical error rate decreases the probability of most cosets, since the trivial (no-error) coset becomes increasingly dominant. 
This reduction in coset probability shows the same quantitative behavior as the observed anomaly and the shift of its critical point.  
Consequently, the probabilities involved in the decoding algorithm become progressively smaller, suggesting numerical underflow as a plausible explanation for the observed anomaly.
Since numerical underflow is itself a deterministic effect, this is consistent with the reproducibility of the anomalous logical error-rate curves.

It is worth noting that the anomaly appears to occur well below the threshold for this noise model.
However, without understanding the underlying mechanism, it it not possible to determine whether the anomaly already affects the results at higher physical noise levels.
As it could happen at small enough scale not being visible, but still influencing the results. 
The next section therefore introduces error mitigation and detection approaches.

\section{Mitigation Strategies}\label{sec:num_underflow_solutions}
In the following we will present three different approaches to reduce the occurrence of the numerical underflow, 
or to at least to reduce the impact it has on the logical error rate curves.

Of the three following approaches we were only able to implement the logarithm and rejection rate, 
either due to time constraints or due to limitation of the used implementation, which could themselves be avoided with more time. 

\subsection{Precision}
% increase precision  -> use better suited programming language
Given that we assume that we are dealing with a numerical underflow problem an obvious problem mitigation strategy is to increase the precision. 
Unfortunately, some parts of our simulations\footnote{
    The problems lies in the choice of programming language. 
    We are using \texttt{Python}, utilizing \texttt{Numpy} \cite{numpy} and \texttt{Numba} \cite{numba} to just in time compile, for increased performance. 
    Sadly the ability to use \texttt{numpy.float128} offered by \texttt{numpy}, seems to be stopped by parts of \texttt{Numba} and \texttt{Numpy} itself, 
    which does not seem to have implemented the higher precision version. 
    This in combination with general slower performance of Python, due to the missing compilation, motivates an implementation in a different programming language, for example \texttt{C++}.
    Sadly the time for required reimplementation was not there, yet.
} do not support such an increased precision, we therefore are required to take a different approach.

\subsection{Logarithm and Rejection Rate}
% calculated the logarithm of the probability instead of the probability  (arons idea)
We can reduce the susceptibility of our implementation by calculating the logarithm of the coset probability\footnote{
    Aron Marton provided the idea for this solution as he encountered similar problems in his own implementation.
}, instead of the coset probability itself.
We implement this for the probability itself as well as for the normalization of the Gaussian state.
In its calculation we are required to calculate a determinant of a matrix. 
Without going into to much detail, we here take the determinant of sparse matrix, which values become increasingly small with smaller physical probabilities.
Even when implementing the computation of this determinant with a function especially constructed to compute the logarithm of determinant, here we used \texttt{numpy.linalg.slogdet}, 
we can observe a numerical underflow, as the normalization $\Gamma$, which should be greater than $0$, turns negative.

We use a normalization with negative sign as an indication than an numerical underflow occurred.
In those cases we flag the whole decoding process as faulty. 
The shots which contain a decoding event which is flagged as faulty, will be rejected, 
leading to the rejection rate $r_{\textrm{rejection}}$ plotted on all logical error rate curves
\begin{equation}
    r_{\textrm{rejection}} = n_{\textrm{rejection}}/n_{shots}.
\end{equation}
Using this rejection logarithm and rejection rate approach achieves logical error rate curves which scales as predicted, as shown in figure \ref{fig:results_code_capcity_Z_ml_p_log_p_phy_slopes}. 

\subsection{Representative Error}
% change the representative error algorithm -> pre decoding with mwpm 
%   note that arons implementation intrinsicly used this -> thus leading to the less 
One can further reduce the occurrence of the numerical underflow error by using a representative error with higher probability\footnote{
    This possibility became apparent while comparing the implementation of the whole simulation of Aron Marton and the one used in this thesis.
}.
Using the representative error construction algorithm, as cited in \ref{sec:num_underflow_decoder} and introduced in \cite{efficient_decoding_algorithm} , 
leads to high weight representative errors.
As an example an single error event happening on the right edge of the surface code, will be represented by an error of weight $2(d-1)$, as both triggered stabilizers are connected to the left boundary.
Especially for the low physical error rate case one a high distance surface code, such configuration can lead to high weight errors.

If we instead would use a MWPM matching algorithm\footnote{
    Aron Marton recycled the result from the MWPM decoder, as an input for the ML decoder, while our approach was based on the algorithm introduced by Bravyi.
    As such, testing the his implementation of the main algorithm of the decoder in my implementation lead to the same occurrence, 
    while testing it in his implementation lead to an neglectable amount of numerical underflow occurrences.
} to find a representative error the associated representative error might be significantly lower.
This inturn leads to a reduced occurrence the numerical underflow, due to the iterative approach of the algorithm \cite{efficient_decoding_algorithm}. 

The difference in the representative error weight and probability is shown in figure \cref{fig:numerical_underflow_avg_weight,fig:numerical_underflow_avg_prob,fig:numerical_underflow_weight_distr}. 
All figures are based on a distance $d=15$ surface code, with $100k$ syndromes sampled.
The figure \ref{fig:numerical_underflow_avg_weight} shows the average weight of the representative error, for a relevant interval of physical error rates.
We can clearly see the difference in the weights, with the MWPM decoder having in approximately only $50\%$ of the weight of the Bravyi approach. 
The resulting average probability of each approach can be seen in figure \ref{fig:numerical_underflow_avg_prob}, where we observe the associated difference in the probability.
Here the MWPM approach is half of the magnitude of the Bravyi approach. 

This not only hold for the average, as shown in figure \ref{fig:numerical_underflow_weight_distr}, but also for the weight distribution of the representative error over all sampled syndromes. 
One can observe that the distribution are completely disjoint, showing that even the syndrome with the highest weight of the MWPM approach has a lower weight than the lowest weight of the Bravyi approach.

In principle the different choices of representative error should not return a different coset probability.
As such we can use the MWPM approach to find a low weight representative error, thus minimizing the risk of a numerical underflow appearing.
This is a tradeoff between the computational cost of MWPM and the problems associated with numerical underflow. 


\begin{figure}[ht]
    \centering
    \includegraphics[width=\linewidth]{numerical_underflow/numerical_underflow_error_avg_weight.pdf}
    \caption{
        Shown is the average weight of the representative error, for a distance $d=15$ surface code, for a relevant interval of the physical error rates, based on a sample size of $100k$ syndromes.
        }
    \label{fig:numerical_underflow_avg_weight}
\end{figure}



\begin{figure}[ht]
    \centering
    \includegraphics[width=\linewidth]{numerical_underflow/numerical_underflow_error_avg_prob.pdf}
    \caption{
        Shown is the average probability of the representative error, for a distance $d=15$ surface code, for a physical error rate of $p_{phy} = 10\%$, based on a sample size of $100k$ syndromes.
        }
    \label{fig:numerical_underflow_avg_prob}
\end{figure}


\begin{figure}[ht]
    \centering
    \includegraphics[width=\linewidth]{numerical_underflow/numerical_underflow_error_weights_distr.pdf}
    \caption{
            Shown is the representative error weight distribution for the two different approaches, for a distance $d=15$ surface code, based on a sample size of $100k$ syndromes.
        }
    \label{fig:numerical_underflow_weight_distr}
\end{figure}





